\documentclass[10pt,a4paper]{amsart}
\usepackage[margin=19mm]{geometry}
\usepackage{amsmath,amssymb,mathtools}
\usepackage[scr=rsfs,cal=euler]{mathalfa}
\usepackage{xfrac}
\usepackage[unicode,hidelinks,bookmarks=false]{hyperref}
\newcommand{\ie}{i.e.\ }
\newcommand{\cf}{cf.\ }
\newcommand{\resp}{respectively\ }
\newcommand{\Subsection}{\subsection}
\providecommand{\nobreakdash}{\nobreak\hbox{-}}
\setlength{\emergencystretch}{2em}
\multlinegap0pt
\usepackage{bm}
\usepackage{bbm}
\usepackage{dsfont}
\usepackage{xspace}
\usepackage{cancel}
\usepackage{mathtools}
\usepackage{amsthm}
\makeatletter
\@ifundefined{thm}{\newtheorem{thm}{Thm}}{}
\@ifundefined{lem}{\newtheorem{lem}[thm]{Lemma}}{}
\makeatother
\def\F#1{f(#1)}
\title[Product-free subsets of $(0,1)$]{Product-free subsets of $(0,1)$}
\author{Leonardo Franchi, W. T. Gowers, Fredy Yip}
\date{Source: arXiv:2607.06073v1 (CC BY 4.0)}
\begin{document}
\maketitle

\noindent\emph{Source and licence.} Product-free subsets of $(0,1)$. Version arXiv:2607.06073v1 (CC BY 4.0), distributed under the Creative Commons Attribution 4.0 International licence, as recorded in the arXiv licence metadata for this exact version and shown on the abstract page https://arxiv.org/abs/2607.06073v1. Full rights notice: licenses/arxiv-2607.06073-CC-BY-4.0.txt.\par\smallskip

\noindent\textit{Editorial scope.} Counted unit: the Lem whose source label is \texttt{double boost lem} in the article (source0.tex lines 2464--2474, source label \texttt{double boost lem}), delivered together with the article's own complete original proof of it (source0.tex lines 2475--2627, one \texttt{proof} environment of 153 lines). No mathematics is written, repaired or re-derived: statement and proof are the article's own text line for line. Required context retained uncounted: minterv (lines 882--904); small case (lines 1176--1178). Every definition, display and label of those lines is retained unchanged. Omitted from this package and not redistributed: 1--881, 905--1175, 1179--2463, 2628--3047. Nothing inside a retained excerpt is omitted or altered: every retained body is delivered exactly as the article's lines are written, so no omission receipt of any kind accompanies this package. Citations: the retained excerpts carry no citation command at all, so no bibliography entry of the article is transcribed and no reference is redistributed. Reference adaptation: none was needed and none was made; every reference inside the delivered unit resolves inside it, so no unresolved reference or citation is produced. Printed slips kept literal: no printed word, symbol, hypothesis or number of the counted statement or of its proof was altered. Layout and class: the article's own class is \texttt{amsart} with packages including amsmath, wrapfig, dsfont, a4wide, amsthm, amssymb, amsfonts, graphicx, fancyhdr, xspace, psfrag, setspace, supertabular, color; the wrapper is the lane's pinned \texttt{amsart} editorial wrapper at 10pt on A4, which restates the theorem-like environments the retained text needs and adds the article's own macro definitions that the retained text needs (\texttt{\textbackslash F}), with extra wrapper packages (bm, bbm, dsfont, xspace, cancel, mathtools). The delivered numbering is the unit's own. Assets: no retained span contains a figure environment, a graphics-inclusion command, a PDF-page inclusion, a table or a TikZ picture; no image file is copied into this package, the package declares no asset dependency and no image file is redistributed with it. Licence: version arXiv:2607.06073v1 of this article is distributed under the Creative Commons Attribution 4.0 International licence, as recorded in the arXiv licence metadata for this exact version and shown on the abstract page https://arxiv.org/abs/2607.06073v1. A full rights notice is delivered at licenses/arxiv-2607.06073-CC-BY-4.0.txt. Characters: every retained excerpt is pure ASCII, so the delivered engine, XeLaTeX with its default Latin Modern text font, reports no missing character.
\begin{lem}\label{minterv}
Let $A$ be a finite union of finite closed intervals, let $I=[a,b]$ be discrepancy maximizing with respect to $A\cap[0,x+1]$,
and let $[\alpha,\beta]\subseteq[0,x+1]$. Then
$$
    |(A-A)\cap [b-\beta-m_A^-(\beta),\, b-\alpha]|
    \geq 2|A\cap[\alpha,\beta]|=2(\F{\beta}-\F{\alpha}).
$$
The corresponding sumset estimate is
$$
    |(A+A)\cap [a+\alpha,\, a+\beta+m_A^-(\beta)]|
    \geq 2|A\cap[\alpha,\beta]|=2(\F{\beta}-\F{\alpha}).
$$
In particular, if $m_A^-(\beta)=0$, then
$$
    |(A-A)\cap [b-\beta,\, b-\alpha]|
    \geq 2|A\cap[\alpha,\beta]|
$$
and
$$
    |(A+A)\cap [a+\alpha,\, a+\beta]|
    \geq 2|A\cap[\alpha,\beta]|.
$$
\end{lem}

\begin{lem}
\label{small case} If $y<3+s$, then $a\geq v$. 
\end{lem}

\begin{lem} \label{double boost lem}
    Assume that $u>2$. Then
    \begin{align*}
        |(A - A)\cap [1, b - 1]|
        &\geq 2\F{b - 1}
        + |A\cap [a, b - m_A^-(u)]|
        + |A\cap [a, b - m_A^-(u'')]|\\
        &\geq 2\F{b - 1} + s + t - m_A^-(u) - m_A^-(u'')\\
        &\geq 2\F{b - 1} + s + t - \F{u} .
    \end{align*}
\end{lem}

\begin{proof}
     Since we are in the remaining case $y<3+s$, Lemma \ref{small case} gives $a\geq v$. Since $v-u\geq1$, we get
    $$
        b-1=a+s-1\geq v+s-1\geq u+s.
    $$
    Moreover, $u>2$ implies $v''\leq u$. Hence, by the previous lemma,
    $$
        u-u''\geq v''-u''\geq s+g\geq s.
    $$
    Also,
    $$
        \F{u+s}=\F{u}
        \qquad\text{and}\qquad
        \F{u''+s}=\F{u''},
    $$
    because $A$ has no mass in the gaps $(u,u+s)$ and $(u'',u''+s)$.
Since $m_A^-(b-1)=0$, we can apply Lemma \ref{minterv}
to the interval $[u+s,b-1]$ to get
$$
    |(A-A)\cap[1,b-u-s]|
    \geq 2|A\cap[u+s,b-1]|
    =
    2(\F{b-1}-\F{u+s}).
$$
Since $u\in A$, the translate by $u$ gives
    $$
        (A\cap[a,b-m_A^-(u)])-u
        \subseteq
        (A-A)\cap[b-u-s,b-u-m_A^-(u)],
    $$
    and therefore contributes
    $$
        |A\cap[a,b-m_A^-(u)]|.
    $$
    Applying Lemma \ref{minterv} to the interval $[u''+s,u]$ we get
$$
    |(A-A)\cap[b-u-m_A^-(u),b-u''-s]|
    \geq 2|A\cap[u''+s,u]|
    =
    2(\F{u}-\F{u''+s}).
$$
Similarly, since $u''\in A$, the translate by $u''$ gives
    $$
        (A\cap[a,b-m_A^-(u'')])-u''
        \subseteq
        (A-A)\cap[b-u''-s,b-u''-m_A^-(u'')],
    $$
    and therefore contributes
    $$
        |A\cap[a,b-m_A^-(u'')]|.
    $$
        Finally, applying Lemma \ref{minterv} to the interval $[1,u'']$ we get
    $$
        |(A-A)\cap[b-u''-m_A^-(u''),b-1]|
        \geq 2|A\cap[1,u'']|
        =
        2\F{u''}.
    $$
    The intervals just constructed are disjoint, except possibly at endpoints. Indeed, using
    $m_A^-(u),m_A^-(u'')\leq t\leq s$ and $u-u''\geq s$, their endpoints are ordered as
    $$
        1\leq b-u-s\leq b-u-m_A^-(u)
        \leq b-u''-s\leq b-u''-m_A^-(u'')\leq b-1.
    $$
    Summing the contributions gives
    \begin{align*}
        |(A-A)\cap[1,b-1]|
        &\geq
        2(\F{b-1}-\F{u+s})
        + |A\cap[a,b-m_A^-(u)]|\\
        &\quad
        +2(\F{u}-\F{u''+s})
        + |A\cap[a,b-m_A^-(u'')]|
        +2\F{u''}.
    \end{align*}
    Since $\F{u+s}=\F{u}$ and $\F{u''+s}=\F{u''}$, we obtain
    $$
        |(A-A)\cap[1,b-1]|
        \geq
        2\F{b-1}
        + |A\cap[a,b-m_A^-(u)]|
        + |A\cap[a,b-m_A^-(u'')]|.
    $$

    For every $z\in\{u,u''\}$, we have
    $$
        |A\cap[a,b-m_A^-(z)]|
        \geq
        |A\cap[a,b]|-m_A^-(z)
        =
        \frac{s+t}{2}-m_A^-(z).
    $$
    Hence,
    $$
        |(A-A)\cap[1,b-1]|
        \geq
        2\F{b-1}+s+t-m_A^-(u)-m_A^-(u'').
    $$

    It remains to prove that
    $$
        m_A^-(u)+m_A^-(u'')\leq \F{u}.
    $$
    First, for every $r\leq u''$ we have
    $$
        2|A\cap[r,u'']|-(u''-r)\leq \F{u''},
    $$
    and therefore
    $$
        m_A^-(u'')\leq \F{u''}.
    $$
    We claim that
    $$
        m_A^-(u)\leq \F{u}-\F{u''}.
    $$
    Let $r\leq u$. If $r\geq u''$, then
    $$
        2|A\cap[r,u]|-(u-r)
        \leq |A\cap[r,u]|
        \leq \F{u}-\F{u''},
    $$
    where we used only $|A\cap[r,u]|\leq u-r$.

    Suppose now that $r<u''$. Since $A\cap(u'',v'')=\varnothing$, we have
    \begin{align*}
        2|A\cap[r,u]|-(u-r)
        &=
        \bigl(2|A\cap[r,u'']|-(u''-r)\bigr)\\
        &\quad
        -(v''-u'')
        +\bigl(2|A\cap[v'',u]|-(u-v'')\bigr)\\
        &\leq
        m_A^-(u'')-(v''-u'')+|A\cap[v'',u]|.
    \end{align*}
    Since
    $$
        m_A^-(u'')\leq t\leq s\leq v''-u'',
    $$
    it follows that
    $$
        2|A\cap[r,u]|-(u-r)
        \leq |A\cap[v'',u]|
        \leq \F{u}-\F{u''}.
    $$
    Taking the supremum over $r\leq u$, we get
    $$
        m_A^-(u)\leq \F{u}-\F{u''}.
    $$
which gives
    $$
        m_A^-(u)+m_A^-(u'')\leq \F{u}.
    $$
\end{proof}

\end{document}
